\chapter{Deployment}
\label{ch:deployment}

This chapter covers building, packaging, and deploying a Kiss application to production environments.

\section{Building a WAR File}
\label{sec:deploy-war}

The standard deployment artifact is a WAR\index{WAR file} (Web Application Archive) file:

\begin{lstlisting}[style=shell]
./bld war
\end{lstlisting}

\noindent
This command:

\begin{enumerate}
  \item Compiles all Java source code (core, precompiled, and application classes).
  \item Copies frontend assets, backend scripts, and configuration files into an exploded directory structure.
  \item Copies \texttt{web-secure.xml} into \texttt{web.xml} to enable production CORS restrictions.
  \item Packages everything into \texttt{work/Kiss.war}.
  \item Copies \texttt{web-unsafe.xml} back into \texttt{web.xml} so local development continues with the permissive CORS configuration.
\end{enumerate}

\noindent
The resulting WAR file is a standard Jakarta EE 11 (Servlet 6.1) archive containing:

\begin{lstlisting}[style=ini]
Kiss.war
  index.html              # Frontend entry point
  index.js                # Application bootstrap
  kiss/                   # Framework frontend files
  screens/                # Application screens
  WEB-INF/
    web.xml               # Servlet configuration (secure)
    classes/              # Compiled Java classes
    lib/                  # JAR dependencies
    backend/              # Groovy/Lisp service files
    application.ini       # Application configuration
    log4j2.xml            # Logging configuration
\end{lstlisting}

\section{Web Configuration Variants}
\label{sec:deploy-webxml}

Kiss maintains two source variants of the servlet configuration in \texttt{src/main/core/WEB-INF/}:

\begin{description}[style=nextline]
  \item[\texttt{web-unsafe.xml}] Development configuration.  Allows CORS\index{CORS} from any origin (\texttt{*}), enabling rapid frontend development from any port or device.  This is the default during development.

  \item[\texttt{web-secure.xml}] Production configuration.  Restricts CORS to specific origins (\texttt{localhost:8000} and IntelliJ's debug port 63342), preventing cross-origin attacks from arbitrary domains.  This is automatically used when building the WAR file.
\end{description}

\noindent
The deployed \texttt{web.xml} is not maintained by hand --- it is a copy that \texttt{bld} writes from one of the two variants above based on whether you are building for development or for production.

\begin{warnbox}
Never deploy with the permissive CORS configuration (\texttt{web-unsafe.xml}) in production.  The wildcard CORS policy allows any website to make requests to your backend.  The WAR build automatically substitutes \texttt{web-secure.xml}; update its allowed origins to match your production domain.
\end{warnbox}

\section{Deploying to Tomcat}
\label{sec:deploy-tomcat}

\subsection{Embedded Tomcat}

Kiss ships with an embedded Tomcat\index{Tomcat} 11 in the \texttt{tomcat/} directory.  To deploy:

\begin{enumerate}
  \item Build the WAR: \texttt{./bld war}
  \item Copy to Tomcat: \texttt{cp work/Kiss.war tomcat/webapps/ROOT.war}
  \item Start Tomcat: \texttt{tomcat/bin/startup.sh}
\end{enumerate}

\noindent
The embedded Tomcat listens on the application's back-end port --- 8001 with the default \texttt{portBase} (Section~\ref{sec:build-ports}).  The \texttt{bld} system stamps the port into \texttt{tomcat/conf/server.xml} on every \texttt{setup-tomcat} run.

\begin{notebox}
A production WAR deployed to the embedded Tomcat (\texttt{webapps/ROOT.war} --- \texttt{./bld build} places one there automatically) does not survive the next development run: \texttt{develop} and \texttt{start-backend} remove any stale \texttt{ROOT.war} before starting Tomcat.  If it were left in place, Tomcat would re-explode the WAR over the \texttt{ROOT} directory at startup, replacing the development tree those tasks stage there --- most visibly \texttt{web.xml}, whose production CORS allow-list rejects the cross-origin development front end.  Treat a WAR deployed to the embedded Tomcat as a production-like test, not a persistent installation.
\end{notebox}

\subsection{External Tomcat}

To deploy to a standalone Tomcat 11 installation:

\begin{enumerate}
  \item Build the WAR: \texttt{./bld war}
  \item Copy \texttt{work/Kiss.war} to Tomcat's \texttt{webapps/} directory, typically renaming it to your application's name (for example \texttt{YourApplication.war}).
  \item Restart Tomcat (or wait for auto-deploy if enabled).
  \item Access the application at \texttt{http://hostname:port/}\allowbreak\texttt{YourApplication/}.
\end{enumerate}

\begin{notebox}
The WAR filename becomes the application's URL context path.  Tomcat takes the deployed file name (minus the \texttt{.war} extension) and uses it as the first path segment when serving the application: \texttt{Kiss.war} is reached at \texttt{/Kiss/}, \texttt{Payroll.war} at \texttt{/Payroll/}, and so on.  Kiss builds the WAR as \texttt{Kiss.war} by default, so when copying to a production server you will usually rename it to whatever name you want to appear in the URL --- for example \texttt{cp work/Kiss.war /var/lib/tomcat/webapps/YourApplication.war}.

To deploy as the root application (no path segment in the URL), rename the WAR to \texttt{ROOT.war} or configure the Tomcat context path explicitly.
\end{notebox}

\section{Other Servlet Containers}
\label{sec:deploy-other}

The WAR file is compatible with any Jakarta EE 11 (Servlet 6.1) container:

\begin{itemize}
  \item \textbf{Tomcat 11} --- the reference container for Kiss.
  \item \textbf{Jetty 12} --- lightweight alternative.
  \item \textbf{WildFly} --- full Jakarta EE application server.
  \item \textbf{GlassFish} --- reference Jakarta EE implementation.
\end{itemize}

\noindent
Deployment procedures vary by container but follow the standard WAR deployment process.

\section{Split-System Deployment}
\label{sec:deploy-split}

For high-traffic applications or CDN-based frontends, the frontend and backend can be deployed on separate servers.

\subsection{Frontend Server}

Serve the frontend as static files using any web server:

\begin{itemize}
  \item Nginx, Apache, or a CDN for HTML, CSS, and JavaScript files.
  \item No server-side processing needed --- the frontend is pure static content.
\end{itemize}

\subsection{Backend Server}

Deploy the WAR file (or just the backend portion) to a servlet container.

\subsection{Connecting Them}

Configure the backend URL in the frontend's \texttt{index.js} by setting \texttt{SystemInfo.backendUrl} before the framework initializes:

\begin{lstlisting}[style=javascript]
// In index.js, before Server initialization
SystemInfo.backendUrl = "https://api.myapp.com";
\end{lstlisting}

\noindent
The frontend's \texttt{Server.js} checks for this value and uses it as the base URL for all backend calls.  When not set, it assumes the backend is co-located (same origin).

\begin{notebox}
In a split deployment, update \texttt{web-secure.xml} to include the frontend's origin in the CORS allowed-origins list.
\end{notebox}

\section{Production Configuration Checklist}
\label{sec:deploy-checklist}

Before deploying to production, verify the following settings in \texttt{application.ini}:

\begin{longtable}{P{4.5cm}P{6cm}}
\toprule
\textbf{Setting} & \textbf{Production Value} \\
\midrule
\endhead
\texttt{DatabaseType} & Your database type (PostgreSQL, MySQL, MSSqlServer, Oracle). \\
\texttt{DatabaseHost} & Production database hostname. \\
\texttt{DatabaseName} & Production database name. \\
\texttt{DatabaseUser} & Database username with appropriate permissions. \\
\texttt{DatabasePassword} & Database password. \\
\texttt{MaxWorkerThreads} & Match to expected concurrent request load (default: 30). \\
\texttt{UserInactiveSeconds} & Session timeout in seconds (default: 900 = 15 min). \\
\bottomrule
\end{longtable}

\noindent
Additional steps:

\begin{itemize}
  \item Use \texttt{web-secure.xml} and update allowed CORS origins.
  \item Configure \texttt{log4j2.xml} for production log levels and file rotation.
  \item Ensure the database schema is up to date.
  \item Set JVM memory parameters (\texttt{-Xmx}, \texttt{-Xms}) in the Tomcat startup script.
  \item Store all passwords hashed (SHA-256) in the \texttt{users} table.
  \item Test the application with the production database before going live.
\end{itemize}

\section{Logging in Production}
\label{sec:deploy-logging}

Server logs are written to \texttt{tomcat/logs/catalina.out}.  The default \texttt{log4j2.xml} configuration sets the root logger to \texttt{ERROR} level with console output:

\begin{lstlisting}[style=ini]
# log4j2.xml pattern
%d{yyyy-MM-dd HH:mm:ss} [%-5p] [%-25c{2}] - %m%n
\end{lstlisting}

\noindent
For production, consider:

\begin{itemize}
  \item Setting the root level to \texttt{WARN} or \texttt{ERROR} to reduce log volume.
  \item Adding a \texttt{RollingFileAppender} for log rotation.
  \item Configuring specific loggers for application packages at \texttt{DEBUG} level during troubleshooting.
  \item Using \texttt{Configurator.setLevel()} to change log levels at runtime without a restart (see Chapter~\ref{ch:utilities}).
\end{itemize}

\section{Performance Tuning}
\label{sec:deploy-performance}

\begin{description}[style=nextline]
  \item[Worker threads] The \texttt{MaxWorkerThreads} setting in \texttt{application.ini} controls the size of the service request thread pool.  Set this to match the expected number of concurrent requests.  A typical starting point is 2--4 times the number of CPU cores.

  \item[Connection pool] C3P0 connection pool settings are configured in \texttt{application.ini}.  Key parameters include:
  \begin{itemize}
    \item \texttt{MaxPoolSize} --- maximum database connections (default: varies by config).  A good starting point is 4 times the number of CPU cores.
    \item \texttt{MinPoolSize} --- minimum idle connections to maintain.
    \item \texttt{MaxIdleTime} --- seconds before idle connections are closed.
    \item \texttt{DatabaseTestOnCheckin} --- validate connections when returned to the pool.
  \end{itemize}

  \item[JVM memory] Set \texttt{-Xmx} (maximum heap) and \texttt{-Xms} (initial heap) in the Tomcat startup script or \texttt{CATALINA\_OPTS} environment variable:
  \begin{lstlisting}[style=shell]
export CATALINA_OPTS="-Xms512m -Xmx2g"
  \end{lstlisting}

  \item[Static file caching] In a split deployment, configure the frontend web server (Nginx, Apache) with appropriate cache headers for static assets.  Kiss frontend files rarely change after deployment.
\end{description}

\begin{tipbox}
Monitor \texttt{tomcat/logs/catalina.out} during the first days of production deployment.  Watch for connection pool exhaustion warnings, slow query logs, and memory pressure indicators.
\end{tipbox}
